\chapter{Basic Theory}
\section{Divisibility}

\begin{definition}
    A number $n$ is divisible by $m$ if there exists some integer $r$ such that $mr=n$. Also, we say $m$
divides $n$ and write this symbolically as
\[
m \mid n
\] 
\end{definition}

\begin{theorem}
    Let $x$, $y$, $z$ be integers.

    \begin{itemize}
        \item We have $x \mid x$.
        \item We always have $1 \mid x$ and $x \mid 0$.
        \item If $x \mid y$ and $y \mid z$, then $x \mid z$.
    \end{itemize}
\end{theorem}


\begin{theorem}

If $ x, y \in \mathbb{Z}$, and $x \mid y$ then $y=0$ or $\mid x \mid \leq \mid y \mid$

\end{theorem}


\begin{theorem}

Suppose $ c, x, y \in \mathbb{Z}$, and $c \mid x, y$ then $c \mid a x + b y$ for any $a, b \in \mathbb{Z}$
 
\end{theorem}

\section{Euclid's Division Lemma}

\begin{theorem}[Euclid's Division Lemma]
    For any integers $b$, $a$, we can find \textbf{unique}
integers $q$, $r$ such that

\[b = aq + r, 0 \leq r < a.\]

Namely, $q$ is called the \textbf{quotient} and $r$ the \textbf{remainder}.
\end{theorem}


\section{Primes}

\begin{definition}
    A natural number is called a \textbf{prime} if it has only two positive divisors, $1$ and itself.
\end{definition}

\begin{theorem}[Euclid]
    There are infinitely many primes.
\end{theorem}

\section{Fundamental Theorem of Arithmetic}

\begin{theorem}[Fundamental Theorem of Arithmetic]
    Any natural number greater than $1$ has a \textbf{unique} prime factorization upto order.Thus, any number $n$ can be written as 
    \[n=p_1^{\alpha_1}p_2^{\alpha_2}\dots p_k^{\alpha_k}\]
    Where $p_1, p_2, \dots, p_k$ are primes and $\alpha_1, \alpha_2,\dots, \alpha_k$ are non-negative integers. 
\end{theorem}

\begin{lemma}
    $\sqrt{p}$ is irrational for any prime $p$.
\end{lemma}

\section{GCD and LCM}

\begin{definition}
    The GCD, or the \textbf{Greatest Common Divisor} of two numbers is the
number obtained by the set of common prime factors. For two numbers $m$, $n$, it is denoted
by $\gcd{m}{n}$, or often just $(m, n)$
\end{definition}

\begin{lemma}
    Let $a$, $b$ be integers. The \textbf{GCD} of $a$, $b$ is the largest number which divides both $a$, $b$. In particular, $\gcd{a}{b} \leq a, b$.
\end{lemma}

\begin{lemma}

Let $a,b,c$ be three integers. Then

\[c\mid a, c\mid b,\Rightarrow  c \mid \gcd(a,b)\]

\end{lemma}

\begin{theorem}[Prime Factorization of GCD]
    Let $a$, $b$ be two integers with prime factorization:
    \[a=p_1^{\alpha_1}p_2^{\alpha_2}\dots p_k^{\alpha_k}\]
    \[b=p_1^{\beta_1}p_2^{\beta_2}\dots p_k^{\beta_k}\]
    where $\alpha_i$, $\beta_i$ are non-negative integers. then
    
    \[\gcd{a}{b}=p_1^{\min{\alpha_1}{\beta_1}}p_2^{\min{\alpha_2}{\beta_2}}\dots p_k^{\min{\alpha_k}{\beta_k}}\]
\end{theorem}

\begin{definition}
    The LCM, or the \textbf{Least Common Multiple} of two numbers is the \textbf{least} natural divisible by \textbf{both} a,b. It is denoted
by $\lcm{a}{b}$, and in particular $a$, $b$ $\leq \lcm{a}{b}$.
\end{definition}

\begin{definition}[LCM]
    Let $a$, $b$ be two integers with prime factorization:
    \[a=p_1^{\alpha_1}p_2^{\alpha_2}\dots p_k^{\alpha_k}\]
    \[b=p_1^{\beta_1}p_2^{\beta_2}\dots p_k^{\beta_k}\]
    where $\alpha_i$, $\beta_i$ are non-negative integers. then
    
    \[\lcm{a}{b}=p_1^{\max{\alpha_1}{\beta_1}}p_2^{\max{\alpha_2}{\beta_2}}\dots p_k^{\max{\alpha_k}{\beta_k}}\]
\end{definition}

\begin{lemma}

Let $a,b,c$ be three integers. Then
\[a\mid c, b\mid c,\Rightarrow \lcm(a,b) \mid c
\]

\end{lemma}

\begin{theorem}[Product of GCD and LCM]. 
    Let $a$, $b$ be two integers. Then
        $\gcd(a, b) \lcm(a, b) = ab.$
\end{theorem}

\begin{definition}
    Two numbers $a$, $b$ are called coprime or relatively prime if $\gcd(a, b) = 1$.
\end{definition}

\begin{lemma}[Euclid's Lemma]
    If $c \mid ab$ and $\gcd{c}{a}=1$, then $c \mid b$.
\end{lemma}

\begin{lemma}
    Let $a$, $m$, $n$ be positive integers. Then,

    \[\gcd{a^{m} - 1}{a^{n} - 1} = a^{\gcd{m}{n}} - 1.\]

\end{lemma}


\begin{lemma}
    If $p$ is an odd prime, and $a, b$ are coprime, then

    \[ \gcd{\frac{a^p+b^p}{a+b}}{a+b} \in \{1,p\}\]
\end{lemma}

\section{Euclid's Division Algorithm}

\begin{lemma}
    Let $a$, $b$ be integers. We can write $a = bq+r$ for integers $q$, $r$ where $0 \leq r < b$.
Then the lemma states that
\[gcd(a, b) = gcd(r, b).\]

\end{lemma}

\section{Bézouts' Theorem}

\begin{theorem}[Bézouts' Theorem]
Let $a$, $b$ be integers. Then the equation

\[ax+by=n\]

has a solution if and only if $\gcd{a}{b}$ divides $n$.
\end{theorem}

\section{Chicken Mc Nuggets}


\begin{theorem}[Chicken Mc Nuggets Theorem/ Frobenius Coin Problem]
Let $a$, $b$ be coprime integers. Then the largest number that can not be expressed on the form 

\[am+bn, \ m,n \in \mathbb{Z}_{\geq 0}\]

Is $(m-1)(n-1)-1$.
\end{theorem}


\section{Base Systems}
\begin{definition}
    Given a positive integer $b$, the base$-b$ representation of a real number $r$ is a string of digits:
    
    \[
    {(a_n a_{n-1} \dots a_1 a_0.a_{-1} a_{-2} \dots)}_{b}
    \]

    Where $a_i \in \{0,1,\dots, b-1\}$, and

    \[
    r= a_n b^n+a_{n-1} b^{n-1} \dots a_1 b^{1}+ a_0 b^{0}+a_{-1}b^{-1} a_{-2}b^{-2} \dots
    \]
\end{definition}

\begin{theorem}
    The base$-b$ representation of any real number is unique.
\end{theorem}


\section{Four Number Lemma}

\begin{theorem}[Four number Lemma]
Let $a$, $b$, $c$, and $d$ be positive integers such that $ab = cd$. Show that there exists positive
integers $p$, $q$, $r$, $s$ such that
\[a = pq, b = rs, c = ps, d = qr\]
\end{theorem}



